% Provenance (never printed): see sources/outline_sources.md (Year 3:
% partition coefficients; batch and fractional melting; crystallisation
% models; reading a spider diagram quantitatively). Brief:
% sources/book5/BRIEFS.md, ch. 12. Curves from figdata/bachelor-3/
% melting_models.py. Exercise and problem values in the Book 5 scratch
% calculator calc/ch12.py. Aldren Island and its lavas are invented.
\chapter{Trace-Element Modelling}\label{ch:b3:trace-element-modelling}

Two basalts erupted on the same island a few kilometres apart. One holds
over one and a half times as much lanthanum as the other, and nearly four times
less ytterbium. Did they melt from two different mantles, or from one mantle
melted by different amounts, or at different depths? The major elements
barely tell: both are basalts. The trace elements, which follow the minerals
without changing them, keep a record of how much melt was made, how it left
its source and which minerals stayed behind. This chapter writes that record
as equations and reads it back.

\begin{recall}
The mineral--melt partition coefficient $D$, the bulk partition coefficient,
compatible and incompatible elements, the rare earth elements, the CI
chondrite and the primitive mantle, normalised diagrams and the europium
anomaly, MORB and OIB are in the Year 2 volume, in its chapter on elements in
the Earth. Peridotite, lherzolite, partial melting and the melt fraction $F$
are in its chapter on melting the mantle; equilibrium and fractional
crystallisation, the Rayleigh fractionation law, mixing and
assimilation--fractional crystallisation (AFC), in its chapter on magmatic
differentiation. There the constancy of the ratio of two very incompatible
elements during crystallisation was proved.
\end{recall}

\section{The bookkeeping of a melting rock}\label{sec:b3:trace-element-modelling:balance}

A rock of mass 1 melts. At the melt fraction $F$ it has made a mass $F$ of
melt and kept a mass $1 - F$ of solid residue. A trace element present at the
concentration $C_0$ in the source is shared between them; its concentration
in the melt is $C_L$, in the residue $C_S$. Every model of this chapter is a
form of one balance, $C_0 = FC_L + (1 - F)C_S$ for the whole melt and residue,
and a rule that links $C_S$ to $C_L$: at the contact between a crystal and
the melt, $C_{\text{mineral}} = D\,C_L$. The models differ in what touches
what, and when.

The partition coefficients of the minerals of the mantle differ widely. In
one compilation, lanthanum has $D = 0.058$ in clinopyroxene and $0.001$ in
garnet and in olivine, orthopyroxene and spinel; ytterbium has $D = 0.58$ in
clinopyroxene, $4.5$ in garnet and $0.16$ in the other three. Garnet takes the
heavy rare earths and rejects the light ones.
% ledger: budahn91:d-ree

\begin{definition}[Non-modal melting]\label{def:b3:trace-element-modelling:nonmodal}
In \emph{non-modal melting}\index{non-modal melting}, the minerals of the
source enter the melt in proportions $p_i$ that differ from their mass
fractions $x_i$ in the source. The melt then sees the bulk partition
coefficient of the source, $D_0 = \sum_i x_iD_i$, at the start and the
melting-weighted coefficient $P = \sum_i p_iD_i$ as it grows.
\end{definition}

Melting is seldom modal. In one worked example of a spinel lherzolite of
\qty{58}{\percent} olivine, \qty{27}{\percent} orthopyroxene,
\qty{12}{\percent} clinopyroxene and \qty{3}{\percent} spinel, the minerals
enter the melt in the proportions 20, 25, 45 and \qty{10}{\percent}:
clinopyroxene melts four times faster than its share, olivine three times
slower.
% ledger: white13:ex73

\begin{proposition}[Bulk partition coefficient of the residue]
\label{prop:b3:trace-element-modelling:residue}
In \omterm{def:b3:trace-element-modelling:nonmodal}{non-modal melting} at constant $p_i$, the residue at the melt fraction $F$
has the bulk partition coefficient
\[
\bar D(F) = \frac{D_0 - PF}{1 - F},
\]
valid while every mineral remains, that is for $F < x_i/p_i$ for each $i$.
\end{proposition}

\begin{proof}
Per unit mass of source, the mass of mineral $i$ left at $F$ is $x_i - p_iF$,
and the residue has mass $1 - F$. Its bulk coefficient is the mass-weighted
mean, $\sum_i (x_i - p_iF)D_i/(1 - F) = (D_0 - PF)/(1 - F)$. A mineral is
exhausted when $x_i - p_iF = 0$.
\end{proof}

With the mineral coefficients above, a spinel lherzolite of the worked
example, and a garnet lherzolite held at \qty{5}{GPa} and
\qty{1400}{\celsius} in the laboratory, of \qty{54.1}{\percent} olivine,
\qty{9.8}{\percent} orthopyroxene, \qty{19.5}{\percent} clinopyroxene and
\qty{16.7}{\percent} garnet by mass, have the bulk coefficients below.
% ledger: white13:ex73, rfr25:gkr-5gpa, budahn91:d-ree
\begin{center}
\small
\begin{tabular}{lccc}
\hline
source & $D_{\mathrm{La}}$ & $D_{\mathrm{Sm}}$ & $D_{\mathrm{Yb}}$ \\
\hline
spinel lherzolite & 0.0078 & 0.062 & 0.21 \\
garnet lherzolite & 0.012 & 0.107 & 0.97 \\
\hline
\end{tabular}
\end{center}
Lanthanum is very incompatible in both. Ytterbium is moderately incompatible
in the spinel lherzolite and nearly neutral in the garnet lherzolite. The
difference between the two sources lies almost entirely in the heavy rare
earths.

\section{Batch melting}\label{sec:b3:trace-element-modelling:batch}

\begin{theorem}[Batch melting]\label{thm:b3:trace-element-modelling:batch}
In \emph{batch melting}\index{batch melting}, the melt stays in equilibrium
with the whole residue until it leaves in one batch. With a constant bulk
partition coefficient $D$, the melt and the residue at the melt fraction $F$
have
\[
\frac{C_L}{C_0} = \frac{1}{D + F(1 - D)}, \qquad
\frac{C_S}{C_0} = \frac{D}{D + F(1 - D)}.
\]
\end{theorem}

\begin{proof}
At equilibrium every crystal of the residue holds $D_iC_L$, so the residue
holds $C_S = DC_L$ on average. The balance $C_0 = FC_L + (1 - F)DC_L$ gives
$C_L$, and $C_S$ follows.
\end{proof}

The formula is the equilibrium crystallisation law of the Year 2 volume read
backwards: in both, $F$ is the fraction of liquid in a closed system. Its two
limits are worth knowing. As $F \to 0$, $C_L \to C_0/D$: the first drop of melt
can enrich an element at most by the factor $1/D$. As $D \to 0$, $C_L \to
C_0/F$: a perfectly incompatible element is diluted into the melt, and its
enrichment measures $F$ alone. For a compatible element, $D > 1$, the melt is
poorer than the source, and the more so the smaller $F$.

\begin{proposition}[Non-modal batch melting]\label{prop:b3:trace-element-modelling:nonmodal}
In \omterm{thm:b3:trace-element-modelling:batch}{batch melting} with the melting proportions $p_i$ of \cref{def:b3:trace-element-modelling:nonmodal},
\[
\frac{C_L}{C_0} = \frac{1}{D_0 + F(1 - P)}.
\]
\end{proposition}

\begin{proof}
The residue holds $\bar D(F)\,C_L$ (\cref{prop:b3:trace-element-modelling:residue}),
so $C_0 = FC_L + (1 - F)\bar D(F)C_L = C_L(F + D_0 - PF)$.
\end{proof}

When $P$ exceeds $D_0$, as for lanthanum in the worked example, where $D_0 =
0.0074$ and $P = 0.025$, the residue loses its best host, clinopyroxene, faster
than the rest. The melt is then a little richer than modal melting predicts,
by under \qty{2}{\percent} at $F = 0.1$ for lanthanum (exercise~6), but the
residue runs out of clinopyroxene at $F = 0.12/0.45 \approx 0.27$, and beyond
that the minerals and coefficients change.
% ledger: white13:ex73

\section{Fractional and dynamic melting}\label{sec:b3:trace-element-modelling:fractional}

Melt in a peridotite gathers along grain edges and can escape as soon as it is
connected. The opposite end member to \omterm{thm:b3:trace-element-modelling:batch}{batch melting} removes every drop as it
forms.

\begin{theorem}[Fractional melting]\label{thm:b3:trace-element-modelling:fractional}
In \emph{fractional melting}\index{fractional melting}, each increment of melt
is in equilibrium with the residue only at the moment it forms, and then leaves
it. With a constant $D$, after a total mass $F$ of melt has been extracted, the
residue and the melt forming at that moment have
\[
\frac{C_S}{C_0} = (1 - F)^{1/D - 1}, \qquad
\frac{C_L}{C_0} = \frac{1}{D}(1 - F)^{1/D - 1}.
\]
\end{theorem}

\begin{proof}
The residue has mass $1 - F$ and concentration $C_S$. When a mass
$\mathrm dF$ of melt leaves at $C_L = C_S/D$, the element in the residue
falls by $\mathrm d[(1 - F)C_S] = -(C_S/D)\,\mathrm dF$. Expanding,
$(1 - F)\,\mathrm dC_S = C_S(1 - 1/D)\,\mathrm dF$, so
$\mathrm dC_S/C_S = (1/D - 1)\,\mathrm d[-\ln(1 - F)]$. Integrating from
$C_S = C_0$ at $F = 0$ gives $\ln(C_S/C_0) = (1/D - 1)\ln(1 - F)$, and
$C_L = C_S/D$.
\end{proof}

For an incompatible element, $1/D - 1$ is large and positive: the residue is
stripped within a few times $D$ of melting, and the instantaneous melt falls
from $C_0/D$ to almost nothing. A residue of \omterm{thm:b3:trace-element-modelling:fractional}{fractional melting} is far more
depleted than one of \omterm{thm:b3:trace-element-modelling:batch}{batch melting}. The melts are not erupted one by one,
though: they gather in channels and pool, and the pooled melt, defined below,
is drawn dashed in the figure.

\begin{omfigure}
\begin{tikzpicture}
\begin{groupplot}[group style={group size=2 by 1, horizontal sep=1.6cm},
  omrecord, width=0.47\linewidth, height=0.44\linewidth, title style={font=\small},
  xlabel={melt fraction $F$}, xmin=0, xmax=0.3, ymode=log, ymin=0.01, ymax=150,
  xtick={0,0.1,0.2,0.3}, xticklabels={0,0.1,0.2,0.3},
  ytick={0.01,0.1,1,10,100}, yticklabels={0.01,0.1,1,10,100},
  legend style={font=\scriptsize, legend cell align=left}]
\nextgroupplot[title={(a) batch (solid), aggregated (dashed)},
  ylabel={$C_L/C_0$}, legend style={at={(0.97,0.97)}, anchor=north east}]
% computed in figdata/bachelor-3/melting_models.py
\addplot[omDef, thick] table[x=f, y=b001] {figdata/out/bachelor-3-melting_models-curves.dat};
\addlegendentry{$D = 0.01$}
\addplot[red!70!black, thick] table[x=f, y=b01] {figdata/out/bachelor-3-melting_models-curves.dat};
\addlegendentry{$D = 0.1$}
\addplot[black, thick] table[x=f, y=b1] {figdata/out/bachelor-3-melting_models-curves.dat};
\addlegendentry{$D = 1$}
\addplot[omLithLime!70!black, thick] table[x=f, y=b10] {figdata/out/bachelor-3-melting_models-curves.dat};
\addlegendentry{$D = 10$}
\addplot[omDef, thick, densely dashed] table[x=f, y=a001] {figdata/out/bachelor-3-melting_models-curves.dat};
\addplot[red!70!black, thick, densely dashed] table[x=f, y=a01] {figdata/out/bachelor-3-melting_models-curves.dat};
\addplot[omLithLime!70!black, thick, densely dashed] table[x=f, y=a10] {figdata/out/bachelor-3-melting_models-curves.dat};
\nextgroupplot[title={(b) fractional, instantaneous}]
\addplot[omDef, thick] table[x=f, y=i001] {figdata/out/bachelor-3-melting_models-curves.dat};
\addplot[red!70!black, thick] table[x=f, y=i01] {figdata/out/bachelor-3-melting_models-curves.dat};
\addplot[black, thick] table[x=f, y=i1] {figdata/out/bachelor-3-melting_models-curves.dat};
\addplot[omLithLime!70!black, thick] table[x=f, y=i10] {figdata/out/bachelor-3-melting_models-curves.dat};
\end{groupplot}
\end{tikzpicture}

\par\smallskip{\small Enrichment of a trace element in the melt against the
melt fraction, for four bulk partition coefficients. (a) Batch melts (solid)
and \omterm{def:b3:trace-element-modelling:fractional}{aggregated fractional melts} (dashed; for $D = 1$ both stay at 1). (b) The
instantaneous melts of \omterm{thm:b3:trace-element-modelling:fractional}{fractional melting}: an element with $D = 0.01$ is
exhausted from the residue after a few per cent of melting.}\label{fig:b3:trace-element-modelling:curves}
\end{omfigure}

\begin{definition}[Aggregated fractional melt]\label{def:b3:trace-element-modelling:fractional}
The \emph{aggregated fractional melt}\index{aggregated fractional melt} is the
mixture of all the instantaneous melts of a \omterm{thm:b3:trace-element-modelling:fractional}{fractional melting} from $0$ to
$F$, pooled without further reaction with the residue.
\end{definition}

\begin{proposition}[Composition of the aggregated melt]
\label{prop:b3:trace-element-modelling:aggregate}
With a constant $D$, the \omterm{def:b3:trace-element-modelling:fractional}{aggregated fractional melt} has
\[
\frac{\bar C_L}{C_0} = \frac{1 - (1 - F)^{1/D}}{F}.
\]
\end{proposition}

\begin{proof}
The element extracted is the source minus the residue,
$C_0 - (1 - F)C_S = C_0[1 - (1 - F)^{1/D}]$ by \cref{thm:b3:trace-element-modelling:fractional},
spread over a mass $F$ of melt.
\end{proof}

Once $(1 - F)^{1/D}$ is small, the aggregated melt holds the whole of the
element and $\bar C_L \approx C_0/F$, whatever $D$: for very incompatible
elements, pooled fractional melts and batch melts with small $D$ agree, and
both measure $F$.

Real melting lies between the two end members. A residue keeps some melt in
its pores, and the melt that leaves is the excess above that porosity.

\begin{definition}[Dynamic melting]\label{def:b3:trace-element-modelling:dynamic}
In \emph{dynamic melting}\index{dynamic melting} (also called continuous
melting), the residue keeps a constant mass fraction $\phi$ of melt in its
pores, in equilibrium with its crystals, and every melt produced beyond it is
extracted.
\end{definition}

\begin{proposition}[Dynamic melting]\label{prop:b3:trace-element-modelling:dynamic}
\omterm{def:b3:trace-element-modelling:dynamic}{Dynamic melting} with porosity $\phi$ is \omterm{thm:b3:trace-element-modelling:fractional}{fractional melting} with the effective
partition coefficient $D' = (1 - \phi)D + \phi$, where $F$ is the mass
fraction of melt extracted:
\[
\frac{C_L}{C_0} = \frac{1}{D'}(1 - F)^{1/D' - 1}.
\]
\end{proposition}

\begin{proof}
The residue is crystals with the fraction $1 - \phi$ and pore melt with the
fraction $\phi$, both in equilibrium with the melt at $C_L$, so it holds
$C_S = [(1 - \phi)D + \phi]C_L = D'C_L$. The proof of
\cref{thm:b3:trace-element-modelling:fractional} then applies with $D'$.
\end{proof}

For $\phi = 0$ the model is \omterm{thm:b3:trace-element-modelling:fractional}{fractional melting}; for $\phi \to 1$ it
approaches \omterm{thm:b3:trace-element-modelling:batch}{batch melting}. A porosity of one per cent caps the enrichment of
the first melt at $1/\phi = 100$, however small $D$: the pore melt buffers
the residue.

\begin{omfigure}
\begin{tikzpicture}[font=\small, >=Stealth]
\tikzset{rock/.style={draw, fill=omLithLime!25, minimum width=2.3cm, minimum height=2.0cm},
  melt/.style={draw, fill=red!25}}
% (a) batch
\node[rock] (a) at (0,0) {};
\foreach \x/\y in {-0.7/0.5, 0.3/0.6, -0.2/-0.1, 0.7/-0.4, -0.6/-0.6, 0.4/0.0}
  \fill[red!55] (\x,\y) circle (0.09);
\node[align=center, font=\scriptsize] at (0,-1.45) {residue $1-F$ at $DC_L$\\melt $F$ at $C_L$};
\draw[->, thick, red!70!black] (0,1.05) -- (0,1.7) node[above, font=\scriptsize] {all at once};
\node[font=\small] at (0,-2.25) {(a) batch};
% (b) fractional
\node[rock] (b) at (4,0) {};
\draw[->, thick, red!70!black] (4,1.05) -- (4,1.7);
\foreach \y/\s in {1.85/0.06, 2.05/0.08, 2.27/0.1}
  \fill[red!55] (4,\y) circle (\s);
\node[align=center, font=\scriptsize] at (4,-1.45) {each drop leaves\\as it forms};
\node[font=\small] at (4,-2.25) {(b) fractional};
% (c) dynamic
\node[rock] (c) at (8,0) {};
\foreach \x/\y in {7.4/0.5, 8.4/0.3, 7.8/-0.5, 8.6/-0.6}
  \fill[red!55] (\x,\y) circle (0.06);
\draw[->, thick, red!70!black] (8,1.05) -- (8,1.7);
\foreach \y/\s in {1.85/0.06, 2.05/0.08, 2.27/0.1}
  \fill[red!55] (8,\y) circle (\s);
\node[align=center, font=\scriptsize] at (8,-1.45) {pores keep $\phi$;\\the excess leaves};
\node[font=\small] at (8,-2.25) {(c) dynamic};
\end{tikzpicture}

\par\smallskip{\small Three ways for melt to leave a melting source. (a)
Batch: the melt stays with the whole residue, then leaves. (b) Fractional:
each increment leaves at once; the pooled increments form the aggregated
melt. (c) Dynamic: the residue keeps a small mass fraction $\phi$ of melt in
its pores.}\label{fig:b3:trace-element-modelling:modes}
\end{omfigure}

\section{Melting, then crystallisation}\label{sec:b3:trace-element-modelling:crystallisation}

A melt that leaves its source rarely erupts unchanged. It crystallises
olivine, then pyroxene and plagioclase on its way up, it may assimilate wall
rock, and it may mix with other batches. The laws of the Year 2 volume apply
in turn: a lava is a melt of fraction $F$ of its source, then a liquid that
has kept a fraction $f$ of its mass through fractional crystallisation, so
that, for an element with $D$ in the melting residue and $D_c$ in the
crystallising minerals,
\[
C_{\text{lava}} = C_0\,\frac{1}{D + F(1 - D)}\,f^{\,D_c - 1}.
\]
The second factor is the same for every element that the crystallising
minerals reject, $D_c \approx 0$: it scales them all by $1/f$. Ratios of such
elements are therefore set at the source, and a lava that has lost olivine
keeps the La/Yb of its primary melt while its lanthanum rises. Inverting $F$
from a concentration needs $f$ as well; inverting the source from a ratio does
not. Plagioclase is the exception among the rare earths: it takes europium,
and its removal leaves the europium anomaly of the Year 2 volume.

\section{Reading rare earth patterns}\label{sec:b3:trace-element-modelling:patterns}

A melt of a spinel lherzolite and a melt of a garnet lherzolite share their
light rare earths, which neither residue keeps, and differ in the heavy ones,
which garnet holds. On a chondrite-normalised diagram the garnet melts are
steep and flat-bottomed at the heavy end; the spinel melts are gentler. The
steepness grows as $F$ falls.

\begin{omfigure}
\begin{tikzpicture}
\begin{axis}[omrecord, width=0.8\linewidth, height=0.5\linewidth, ymode=log,
  xmin=56.5, xmax=71.5, ymin=1, ymax=1000, xtick={57,58,60,62,63,64,65,70,71},
  xticklabels={La,Ce,Nd,Sm,Eu,Gd,Tb,Yb,Lu}, ytick={1,10,100,1000}, yticklabels={1,10,100,1000},
  ylabel={sample / CI chondrite},
  legend style={font=\scriptsize, at={(0.98,0.98)}, anchor=north east,
  legend cell align=left, legend columns=2}]
% computed in figdata/bachelor-3/melting_models.py
\addplot[omDef, thick, mark=*, mark size=1.3pt] table[x=z, y=sp01] {figdata/out/bachelor-3-melting_models-ree.dat};
\addlegendentry{spinel, $F = 1\,\%$}
\addplot[red!70!black, thick, mark=square*, mark size=1.3pt] table[x=z, y=gt01] {figdata/out/bachelor-3-melting_models-ree.dat};
\addlegendentry{garnet, $F = 1\,\%$}
\addplot[omDef, thick, densely dashed, mark=*, mark size=1.3pt] table[x=z, y=sp05] {figdata/out/bachelor-3-melting_models-ree.dat};
\addlegendentry{spinel, $5\,\%$}
\addplot[red!70!black, thick, densely dashed, mark=square*, mark size=1.3pt] table[x=z, y=gt05] {figdata/out/bachelor-3-melting_models-ree.dat};
\addlegendentry{garnet, $5\,\%$}
\addplot[omDef, thick, densely dotted, mark=*, mark size=1.3pt] table[x=z, y=sp20] {figdata/out/bachelor-3-melting_models-ree.dat};
\addlegendentry{spinel, $20\,\%$}
\addplot[red!70!black, thick, densely dotted, mark=square*, mark size=1.3pt] table[x=z, y=gt20] {figdata/out/bachelor-3-melting_models-ree.dat};
\addlegendentry{garnet, $20\,\%$}
\addplot[black, thin, mark=none] table[x=z, y=pm] {figdata/out/bachelor-3-melting_models-ree.dat};
\addlegendentry{source}
\end{axis}
\end{tikzpicture}

\par\smallskip{\small Chondrite-normalised rare earth patterns of batch melts
of a primitive-mantle source, from a spinel lherzolite (blue circles) and a
garnet lherzolite (red squares), at three melt fractions; the source is the
thin line. Only the nine rare earths with tabulated partition coefficients are
plotted.}\label{fig:b3:trace-element-modelling:ree}
\end{omfigure}
% ledger: budahn91:d-ree, ms95:pm-trace, lodders:ci-table, white13:ex73, rfr25:gkr-5gpa

Garnet gives way to spinel as the aluminous mineral of the mantle at about 60
to \qty{90}{km} depth. A melt whose heavy rare earths are flat and low, with a
steep light-to-heavy slope, has therefore left garnet in its residue and
began to melt at least that deep.
% ledger: white13:sp-gt

A pattern mixes two unknowns, $F$ and the source. Plotting two elements
against each other separates them.

\begin{definition}[Ratio--ratio diagram]\label{def:b3:trace-element-modelling:ratios}
A \emph{ratio--ratio diagram}\index{ratio--ratio diagram} plots, for a suite
of lavas, the ratio of a highly incompatible element $H$ to a less
incompatible one $M$ against the concentration of $H$ (or against another
ratio). Partial melting, crystallisation and mixing draw lines of different
shapes on it.
\end{definition}

\begin{proposition}[The batch melting line]\label{prop:b3:trace-element-modelling:line}
For batch melts of one source with $D_H = 0$ and a constant $D_M$, the points
$(C_H, C_H/C_M)$ lie on the straight line
\[
\frac{C_H}{C_M} = \frac{D_M}{C_0^M}\,C_H + \frac{C_0^H(1 - D_M)}{C_0^M}.
\]
\end{proposition}

\begin{proof}
With $D_H = 0$, $C_H = C_0^H/F$, so $1/F = C_H/C_0^H$. Then
$C_H/C_M = (C_0^H/F)\,(D_M + F(1 - D_M))/C_0^M = C_0^HD_M/(FC_0^M) +
C_0^H(1 - D_M)/C_0^M$, and the first term is $(D_M/C_0^M)\,C_H$.
\end{proof}

Fractional crystallisation of minerals that reject both elements moves a lava
horizontally on this diagram, since the ratio does not change; melting moves it
along a sloping line whose slope grows with $D_M$; mixing of two melts draws a
hyperbola. A source with garnet has a large $D_{\mathrm{Yb}}$, hence a steep
melting line in La/Yb against La.

\begin{method}[Inverting two lavas]\label{met:b3:trace-element-modelling:invert}
(1) Check that the lavas are cogenetic and that $H$ is far more incompatible
than $M$. (2) Correct each lava for crystallisation, or use only ratios.
(3) Through two or more lavas, fit the line $C_H/C_M = aC_H + b$. (4) Then
$a = D_M/C_0^M$ and $b = C_0^H(1 - D_M)/C_0^M$: with $C_0^M$ assumed, $D_M =
aC_0^M$ and $C_0^H = bC_0^M/(1 - D_M)$. (5) Each lava's melt fraction is
$F = C_0^H/C_H$.
\end{method}

\begin{omfigure}
\begin{tikzpicture}
\begin{axis}[omrecord, width=0.78\linewidth, height=0.52\linewidth,
  xmode=log, ymode=log, xmin=2, xmax=60, ymin=2, ymax=100,
  xlabel={La (\unit{ppm})}, ylabel={La/Yb},
  xtick={2,5,10,20,50}, xticklabels={2,5,10,20,50},
  ytick={2,5,10,20,50,100}, yticklabels={2,5,10,20,50,100},
  legend style={font=\scriptsize, at={(0.03,0.97)}, anchor=north west,
  legend cell align=left}]
% computed in figdata/bachelor-3/melting_models.py
\addplot[omDef, thick] table[x=la_sp, y=lay_sp] {figdata/out/bachelor-3-melting_models-traj.dat};
\addlegendentry{spinel lherzolite}
\addplot[red!70!black, thick] table[x=la_gt, y=lay_gt] {figdata/out/bachelor-3-melting_models-traj.dat};
\addlegendentry{garnet lherzolite}
\addplot[omDef, only marks, mark=|, mark size=3pt, forget plot] table[x=la_sp, y=lay_sp] {figdata/out/bachelor-3-melting_models-ticks.dat};
\addplot[red!70!black, only marks, mark=|, mark size=3pt, forget plot] table[x=la_gt, y=lay_gt] {figdata/out/bachelor-3-melting_models-ticks.dat};
\addplot[black, only marks, mark=diamond*, mark size=2.8pt] table[x=la, y=layb] {figdata/out/bachelor-3-melting_models-lavas.dat};
\addlegendentry{lavas A and B}
\node[font=\scriptsize, anchor=north] at (axis cs:9.6,5.2) {A};
\node[font=\scriptsize, anchor=south east] at (axis cs:15.2,36) {B};
\node[font=\scriptsize, omDef, anchor=north west] at (axis cs:36.5,17.5) {1\,\%};
\node[font=\scriptsize, omDef, anchor=north west] at (axis cs:11.3,6.2) {5\,\%};
\node[font=\scriptsize, omDef, anchor=north west] at (axis cs:3.14,2.55) {20\,\%};
\node[font=\scriptsize, red!70!black, anchor=south east] at (axis cs:29.5,66) {1\,\%};
\node[font=\scriptsize, red!70!black, anchor=south east] at (axis cs:10.5,24) {5\,\%};
\node[font=\scriptsize, red!70!black, anchor=south east] at (axis cs:3.09,7.0) {20\,\%};
\end{axis}
\end{tikzpicture}

\par\smallskip{\small Batch melts of a primitive-mantle source on a La/Yb
against La diagram, for a spinel and a garnet lherzolite, with ticks at $F = 1$,
2, 5, 10 and \qty{20}{\percent}, and the two invented lavas of the weekend
problem. Each lava sits on one curve only.}\label{fig:b3:trace-element-modelling:ratios}
\end{omfigure}

On logarithmic axes the straight line of the proposition bends, and with
$D_{\mathrm{La}} \ne 0$ it is only nearly straight on linear axes; the curves
of the figure are exact. A spinel source cannot make a melt with La/Yb above
$(C_0^{\mathrm{La}}/D_{\mathrm{La}})/(C_0^{\mathrm{Yb}}/D_{\mathrm{Yb}})$,
about 40 for the primitive mantle; a garnet source can reach about 120.

\section{Exercises}

\begin{exercise}[$\star$]\label{exo:b3:trace-element-modelling:1}
With the mineral coefficients of samarium, 0.40 in clinopyroxene, 0.11 in
garnet and 0.016 in olivine, orthopyroxene and spinel, compute the bulk
partition coefficient of samarium in the spinel lherzolite (\qty{12}{\percent}
clinopyroxene) and in the garnet lherzolite (\qty{19.5}{\percent} clinopyroxene,
\qty{16.7}{\percent} garnet) of the chapter.
\end{exercise}
% ledger: budahn91:d-ree

\begin{exercise}[$\star$]\label{exo:b3:trace-element-modelling:2}
An element with $D = 0.05$ is present at \qty{1}{ppm} in a source. Compute
its concentration in batch melts at $F = 0.02$ and $0.10$.
\end{exercise}

\begin{exercise}[$\star$]\label{exo:b3:trace-element-modelling:3}
What is the largest enrichment $C_L/C_0$ that \omterm{thm:b3:trace-element-modelling:batch}{batch melting} can give to an
element with $D = 0.01$? And to an element with $D = 0$ at
$F = \qty{5}{\percent}$?
\end{exercise}

\begin{exercise}[$\star$]\label{exo:b3:trace-element-modelling:4}
A lava holds \qty{30}{ppm} La and \qty{2.5}{ppm} Yb. With the CI chondrite
values La 0.232 and Yb \qty{0.163}{ppm}, compute its La/Yb and its
chondrite-normalised $(\mathrm{La}/\mathrm{Yb})_N$.
\end{exercise}
% ledger: lodders:ci-table

\begin{exercise}[$\star\star$]\label{exo:b3:trace-element-modelling:5}
For an element with $D = 0.01$, compare the residues after \qty{5}{\percent}
of \omterm{thm:b3:trace-element-modelling:batch}{batch melting} and of \omterm{thm:b3:trace-element-modelling:fractional}{fractional melting}. What does the instantaneous
fractional melt hold at that moment?
\end{exercise}

\begin{exercise}[$\star\star$]\label{exo:b3:trace-element-modelling:6}
In the worked example of the chapter (modes 58, 27, 12, \qty{3}{\percent};
melting proportions 20, 25, 45, \qty{10}{\percent} for olivine,
orthopyroxene, clinopyroxene and spinel), lanthanum has the coefficients
0.0003, 0.002, 0.053 and 0.01. Compute $D_0$ and $P$, then the enrichment of
lanthanum in the batch melt at $F = 0.02$ and $0.10$, non-modal and modal.
\end{exercise}
% ledger: white13:ex73

\begin{exercise}[$\star\star$]\label{exo:b3:trace-element-modelling:7}
At $F = 0.10$, compare the batch melt and the \omterm{def:b3:trace-element-modelling:fractional}{aggregated fractional melt} for
$D = 0.1$ and for $D = 0.01$. Which differ more, and why?
\end{exercise}

\begin{exercise}[$\star\star$]\label{exo:b3:trace-element-modelling:8}
An element with $D = 0.001$ melts dynamically with $\phi = 0.01$. Compute
$D'$, the enrichment of the first melt extracted, and that of the melt
extracted at $F = 0.01$; compare with \omterm{thm:b3:trace-element-modelling:fractional}{fractional melting}.
\end{exercise}

\begin{exercise}[$\star\star$]\label{exo:b3:trace-element-modelling:9}
In the spinel lherzolite of the chapter, $D_{\mathrm{La}} = 0.0078$ and
$D_{\mathrm{Sm}} = 0.062$. By what factor does \omterm{thm:b3:trace-element-modelling:batch}{batch melting} raise La/Sm above
the source value at $F = 0.01$ and at $F = 0.10$? What does that say about
using ratios to estimate $F$?
\end{exercise}

\begin{exercise}[$\star\star\star$]\label{exo:b3:trace-element-modelling:10}
Show that the \omterm{def:b3:trace-element-modelling:fractional}{aggregated fractional melt} is the mean of the instantaneous
melts from 0 to $F$, and that it tends to $C_0/F$ when $(1 - F)^{1/D} \to 0$.
For $D = 0.01$, from what $F$ is it within \qty{1}{\percent} of $C_0/F$?
\end{exercise}

\begin{exercise}[$\star\star\star$]\label{exo:b3:trace-element-modelling:11}
Two cogenetic lavas P and Q have La 20 and \qty{10}{ppm} and La/Yb 10 and 6.
Assuming lanthanum perfectly incompatible and a source with
\qty{0.44}{ppm} Yb, apply the method of the chapter: find $D_{\mathrm{Yb}}$,
the La of the source and the melt fraction of each lava.
\end{exercise}

\begin{exercise}[$\star\star\star$]\label{exo:b3:trace-element-modelling:12}
In the \omterm{def:b3:trace-element-modelling:nonmodal}{non-modal melting} of exercise~6, at what melt fraction is each mineral
exhausted? Compute the bulk coefficient of the residue for lanthanum at
$F = 0.2$, and say why the formula of the chapter cannot be used beyond the
first exhaustion.
\end{exercise}

\section{Problem: Two Lavas of Aldren Island}

\begin{problem}[{Weekend problem --- the source, depth and degree of melting
of two invented basalts, from their lanthanum and ytterbium}]\label{pb:b3:trace-element-modelling:1}
On Aldren Island, an invented volcanic island, lava A holds \qty{9.6}{ppm} La
and \qty{1.7}{ppm} Yb, lava B \qty{15.6}{ppm} La and \qty{0.45}{ppm} Yb.
Both are taken as primary melts of a primitive-mantle source with
\qty{0.648}{ppm} La and \qty{0.441}{ppm} Yb (CI chondrite: La 0.232, Yb
\qty{0.163}{ppm}). Mineral coefficients: La 0.058 in clinopyroxene and 0.001
in garnet and in olivine, orthopyroxene and spinel; Yb 0.58 in
clinopyroxene, 4.5 in garnet and 0.16 in olivine, orthopyroxene and spinel.
The spinel lherzolite is \qty{58}{\percent} olivine, \qty{27}{\percent} orthopyroxene, \qty{12}{\percent}
clinopyroxene, \qty{3}{\percent} spinel; the garnet lherzolite \qty{54.1}{\percent} olivine, \qty{9.8}{\percent}
orthopyroxene, \qty{19.5}{\percent} clinopyroxene, \qty{16.7}{\percent} garnet. Melting is modal.
% ledger: ms95:pm-trace, lodders:ci-table, budahn91:d-ree, white13:ex73, rfr25:gkr-5gpa

\medskip\noindent\textbf{Part I --- The two sources.}
\begin{enumerate}
\item Compute $D_{\mathrm{La}}$ and $D_{\mathrm{Yb}}$ of the spinel lherzolite.
\item Compute them for the garnet lherzolite. Round all four to two significant
figures for the rest of the problem.
\item Which coefficient differs most between the sources, and why? Is
ytterbium compatible in the garnet lherzolite?
\item Compute the chondrite-normalised $(\mathrm{La}/\mathrm{Yb})_N$ of the
source. What does its pattern look like?
\end{enumerate}

\medskip\noindent\textbf{Part II --- Batch melting of lava A.}
\begin{enumerate}[resume]
\item Compute La in batch melts of the spinel lherzolite at $F = 1$, 5 and
\qty{10}{\percent}.
\item Find the melt fraction that gives lava A its lanthanum.
\item Predict its ytterbium and its La/Yb, and compare with lava A.
\item Compute La and Yb in the residue, and the fraction of the source's
lanthanum left in it.
\item What melt fraction would you find by taking lanthanum as perfectly
incompatible?
\end{enumerate}

\medskip\noindent\textbf{Part III --- Fractional melting of lava A.}
\begin{enumerate}[resume]
\item Compute La in the instantaneous fractional melt at $F = 0.1$, 1 and
\qty{6}{\percent}.
\item Compute La in the \omterm{def:b3:trace-element-modelling:fractional}{aggregated fractional melt} at $F = \qty{6.0}{\percent}$.
\item Find the melt fraction at which the aggregated melt has the lanthanum of
lava A. Why is it so close to the answer of question~9?
\item Compare the lanthanum of the residue with that of question~8.
\item Which estimate of $F$ do you keep?
\end{enumerate}

\medskip\noindent\textbf{Part IV --- Lava B and garnet.}
\begin{enumerate}[resume]
\item Compute La/Yb and $(\mathrm{La}/\mathrm{Yb})_N$ of both lavas.
\item What is the largest La/Yb that a batch melt of the spinel lherzolite can
reach?
\item At what melt fraction would a spinel-lherzolite melt have the La/Yb of
lava B, and how much lanthanum would it hold? Conclude.
\item Find the melt fraction of lava B from its lanthanum, for the garnet
lherzolite, and predict its ytterbium.
\item What do the two sources say about the depths at which A and B began to
melt? Could lava A come from the garnet lherzolite?
\item Compute the slope and intercept of the \omterm{thm:b3:trace-element-modelling:batch}{batch melting} line La/Yb $= a\,
\mathrm{La} + b$ of the garnet source with $D_{\mathrm{La}}$ taken as zero,
and check lava B against it.
\item If lava A had lost \qty{10}{\percent} of its mass as olivine before
erupting, how would its La/Yb and its melt fraction change?
\item Answer the question of the chapter's opening for these two lavas.
\item Give the degree of melting of lava A.
\end{enumerate}
\end{problem}
